\subsection{Parâmetros Cinemáticos}

A cinemática do manipulador robótico foi modelada com base na convenção clássica de Denavit-Hartenberg, pois essa metodologia permite representar a posição e orientação relativa entre os elos consecutivos por meio de quatro parâmetros:
o deslocamento ao longo do eixo $z_{i-1}$ ($d_i$),
a rotação em torno de $z_{i-1}$ ($\theta_i$),
o comprimento do elo ao longo do eixo $x_i$ ($a_i$) e
a torção entre os eixos $z_{i-1}$ e $z_i$ ($\alpha_i$).

A Tabela \ref{tab:DH_parametros} resume os parâmetros adotados para cada junta $J_i$ do manipulador robótico, sendo que as juntas $J_1$ a $J_4$ são do tipo revolutas, enquanto $J_5$ é prismática e está embutida no interior do Elo 4. Para juntas as revolutas, o parâmetro variável é $\theta_i$, que representa o ângulo de rotação em torno do eixo $z_{i-1}$. Para a junta prismática $J_5$, a variável é $d_5$, que representa o deslocamento linear ao longo do eixo $z_4$, com valores variando entre $0$ e $0{,}15[m]$.

A definição dos eixos coordenados segue os critérios geométricos e de montagem do manipulador robótico, sendo que os eixos $z_i$ coincidem com os eixos de rotação ou translação das juntas, e os eixos $x_i$ são orientados ao longo do elo seguinte.
Por sua vez, o elo 1, por exemplo, estende-se verticalmente, alinhado ao eixo Z da base, enquanto que os elos seguintes seguem a orientação dos respectivos eixos $x_i$.

\begin{table}[ht]
\centering
\caption{Parâmetros de Denavit-Hartenberg do manipulador robótico.}
\label{tab:DH_parametros}
\begin{tabular}{ccccc}
\toprule
Junta $i$ & $\theta_i$ & $d_i$ [m] & $a_i$ [m] & $\alpha_i$ [rad] \\
\midrule
1 & $\theta_1$ & 0{,}005        & 0{,}000  & $\frac{\pi}{2}$ \\ [4pt]
2 & $\theta_2$ & 0{,}100        & 0{,}005  & 0{,}000         \\ [4pt]
3 & $\theta_3$ & 0{,}000        & 0{,}005  & 0{,}000         \\ [4pt]
4 & $\theta_4$ & 0{,}000        & 0{,}005  & 0{,}000         \\ [4pt]
5 & 0          & $d_5 \in [0; 0,15]$ & 0{,}000  & 0{,}000       \\ [4pt]
\bottomrule
\end{tabular}
\end{table}
